\chapter{\nt{GLUT} 뉴런은 무엇을 계산하는가}
\label{sec:glutcomputer}

\section{게임의 규칙}

\ref{sec:neurontypes}장에서 보았듯이 뇌 뉴런의 압도적 다수는 \nt{GLUT}이다. 이들은 흥분만 시키며, 가중치는 양수뿐이다. 이런 뉴런을 원하는 만큼 가져와서 묻자. 살아 있는 몸에서 일어나는 일만 허용한다면, 이런 네트워크는 무엇을 계산할 수 있는가?

살아 있는 몸에는 모든 소자를 한꺼번에 전환하는 클록 발생기가 없다. 뉴런은 저마다 연속 시간에서 독립적으로 작동한다. 스파이크는 제각각의 지연을 두고 아무 때나 오고 간다. 빛, 소리, 압력에 반응하는 감각 뉴런이 \emph{입력}이고, 근육을 제어하는 뉴런이 \emph{출력}이다. 그 사이에 네트워크가 있으며, 누구도 네트워크에 언제 계산을 시작하고 언제 끝낼지 알려 주지 않는다. 전자공학자에게 이것은 소자 지연을 미리 정확히 알 수 없는 \emph{비동기} 회로이다.

뉴런 모델로는 연속 시간에서 제대로 작동하는 모델 가운데 가장 단순한 것을 쓰자. 뉴런 막의 전압 $V$는 휴지 상태에서 0이다. 뉴런 $j$에서 스파이크가 하나 올 때마다 전압이 $w_j\ge0$만큼 뛰어오르고, 그 뒤 시간 상수 $\tau$로 저절로 0까지 흘러내린다.
\begin{empheq}[box=\keybox]{equation}
\tau\,\frac{\dd V}{\dd t} \;=\; -V \;+\; \tau\sum_j w_j \sum_k \delta\bigl(t-t_j^{(k)}-d_j\bigr), \qquad w_j\ge 0 .
\label{eq:glutneuron}
\end{empheq}
여기서 $t_j^{(k)}$는 뉴런 $j$의 스파이크 시각, $d_j$는 그 뉴런에서 오는 경로의 지연, $\delta$는 순간적인 도약을 뜻하는 델타 함수이다. $V$가 문턱값 $\theta$에 이르면 뉴런은 스파이크를 내고, $V$는 0으로 초기화되며, 약 1밀리초 동안 뉴런은 다시 발화하지 못한다. 전형적인 수치는 이렇다. $\tau$는 뉴런에 따라 1밀리초 미만에서 20밀리초까지, 지연은 1밀리초 미만에서 수십 밀리초까지이다. 이것이 \emph{누설 적분 발화 뉴런}, 곧 문턱값이 있는 RC 회로이다. 의도적인 단순화이다. 피질 뉴런에서는 입력 가지가 스스로 국소 스파이크를 내며~\cite{Smith2013}, 뉴런 하나는 오히려 작은 다층 네트워크처럼 행동한다(\ref{sec:synthesis}장). 이 장의 질문에는 RC 회로 하나로 충분하다.

논리 게이트와의 가장 큰 차이는 누설이다. 뉴런은 스파이크를 영원히 기억하지 않고 대략 $\tau$ 동안만 기억하며, 이 창 안에 들어온 것만 더해진다.

\section{OR와 AND}

게이트부터 시작하자. 스파이크 하나가 전압을 문턱값 위로 올리면, 즉 $w\ge\theta$이면 뉴런은 어느 입력에서 오든 모든 스파이크에 발화한다. 이것이 \emph{OR}이다.

AND는 더 흥미롭다. $w<\theta<2w$라고 하자. 스파이크 하나로는 모자라고 둘이 필요하다. 그런데 이제 `둘 다 왔는가'라는 질문은 \emph{어느 시간 안에} 왔는지를 말하기 전에는 의미가 없다. 첫 스파이크가 시각 $0$에, 둘째가 시간 $\Delta$ 뒤에 왔다고 하자. 둘째가 올 때 첫째에서 남은 것은 $we^{-\Delta/\tau}$이고, $we^{-\Delta/\tau}+w\ge\theta$이면 뉴런이 발화한다. 여기서 일치 창이 나온다.
\begin{empheq}[box=\keybox]{equation}
\Delta \;\le\; \Delta_{\max} \;=\; \tau\,\ln\frac{w}{\theta-w} .
\label{eq:window}
\end{empheq}
이것은 논리 AND가 아니라 \emph{동시성 검출기}, 곧 시간 속의 AND이다(그림~\ref{fig:coincidence}). $\theta=1.5w$이면 창은 $\tau\ln2\approx0.7\tau$이다. $\tau=10\ms$인 피질 뉴런에서는 약 7밀리초이다. $\tau$가 1밀리초보다 짧은 청각계 뉴런에서는 창이 1밀리초 미만으로 줄어든다.

\begin{figure}[t]
\centering
\begin{tikzpicture}[>=Stealth,x=0.9cm,y=1.6cm]
\foreach \dx/\lab/\c [count=\i] in {0.5/{일치함}/red!70!black, 2.6/{일치하지 않음}/blue!60!black}{
 \begin{scope}[xshift={(\i-1)*7.2cm}]
  \draw[->,gray!70!black] (0,0) -- (6.3,0) node[below left,black] {\small $t$};
  \draw[->,gray!70!black] (0,0) -- (0,1.75) node[left,black] {\small $V$};
  \draw[densely dashed,gray!60!black] (0,1.5) -- (6.0,1.5) node[right,black] {\small $\theta$};
  \draw[very thick,\c] (0,0) -- (1,0) -- (1,1)
      plot[domain=1:1+\dx,samples=30] (\x,{exp(-(\x-1)/1.5)});
  \pgfmathsetmacro{\v}{exp(-\dx/1.5)+1}
  \draw[very thick,\c] (1+\dx,{exp(-\dx/1.5)}) -- (1+\dx,{min(\v,1.5)});
  \ifdim\v pt<1.5pt
    \draw[very thick,\c] plot[domain=1+\dx:6,samples=40] (\x,{\v*exp(-(\x-1-\dx)/1.5)});
  \else
    \draw[very thick,\c] (1+\dx,1.5) -- (1+\dx,0) -- (6,0);
    \draw[->,thick,\c] (1+\dx,1.55) -- (1+\dx,1.72) node[above,font=\scriptsize] {스파이크};
  \fi
  \draw[->,thick] (1,-0.35) -- (1,-0.08); \draw[->,thick] (1+\dx,-0.35) -- (1+\dx,-0.08);
  \draw[decorate,decoration={brace,amplitude=3pt,mirror}] (1,-0.42) -- (1+\dx,-0.42) node[midway,below=3pt] {\scriptsize $\Delta$};
  \node[\c] at (3.5,-0.95) {\small \lab};
 \end{scope}}
\end{tikzpicture}
\caption{동시성 검출기, 문턱값 $\theta=1.5w$. 왼쪽에서는 둘째 스파이크가 $\Delta<\Delta_{\max}$ 뒤에 와서 합이 문턱값에 이르고 뉴런이 발화했다. 오른쪽은 $\Delta>\Delta_{\max}$이다. 첫 스파이크에서 거의 아무것도 남지 않아 뉴런은 침묵했다.}
\label{fig:coincidence}
\end{figure}

AND를 얻는 다른 방법은 정확한 시각이 아니라 지속 시간을 쓰는 것이다. 빛이 켜져 있는 동안 감각 뉴런은 스파이크를 자주, 고르게 낸다. 그런 입력 하나로는 문턱값에 모자라고 둘이면 충분하다면, 뉴런은 두 입력이 모두 활성인 동안 작동한다. 이렇게 네트워크는 \emph{수준}으로 계산하며, 스파이크가 도착하는 정확한 시각은 중요하지 않다. 아래 내용의 대부분이 바로 이런 방식이다.

\section{할 수 없는 것}

이제 NOT을 만들어 보자. 입력이 조용할 때 작동하는 뉴런이다. 만들 수 없다. 입력 스파이크가 없으면 \eqref{eq:glutneuron}에는 도약이 하나도 없고, 전압은 문턱값 아래인 0에 머문다. 뉴런이 많은 네트워크라면? 역시 안 되며, 이것은 모든 네트워크에 대해 한꺼번에 증명된다.

발화율로 말하는 것이 가장 편하다. $r_i$를 뉴런 $i$의 발화율, $I_i$를 입력에서 오는 유입, $f_i$를 그 뉴런의 전달 특성, 즉 발화율이 총 유입에 어떻게 의존하는지라고 하자. 적분 뉴런의 전달 특성은 감소하지 않는다. 유입이 많을수록 스파이크가 잦다. 평형에 이른 네트워크의 발화율은 다음 방정식을 만족한다.
\begin{equation}
r_i \;=\; f_i\Bigl(\sum_j w_{ij}\,r_j + I_i\Bigr), \qquad w_{ij}\ge 0 .
\label{eq:ratenet}
\end{equation}

\begin{keyblock}
\begin{proposition}[단조성]
\label{prop:monotone}
네트워크~\eqref{eq:ratenet}가 \nt{GLUT} 뉴런만으로 이루어지고 완전한 침묵에서 출발한다고 하자. 입력을 강하게 하면 어떤 뉴런의 정상 상태 발화율도 줄어들지 않는다. 이런 네트워크가 계산하는 것은 모두 입력의 \emph{단조} 함수이다.
\end{proposition}
\end{keyblock}

증명. 침묵 $r^{(0)}=0$에서 출발해 발화율을 \eqref{eq:ratenet}의 우변에 대입해 나가자. $r^{(k+1)}=F\bigl(r^{(k)}\bigr)$. 가중치가 음이 아니고 $f_i$가 감소하지 않으므로 우변 $F$는 어느 인수에 대해서도 감소하지 않는다. $r^{(1)}\ge0=r^{(0)}$이므로 귀납적으로 $r^{(k+1)}\ge r^{(k)}$이다. 발화율은 커지기만 하며 \eqref{eq:ratenet}의 가장 작은 해에 이른다. 입력 $I$를 더 큰 $I'$로 바꾸면 매 단계에서 $r^{(k)}(I)\le r^{(k)}(I')$이고, 극한에서도 그렇다. 실제 네트워크는 단계가 아니라 연속적으로 평형에 이르지만 같은 결론이 성립한다. 연결이 양수이면 두 실행 사이의 부등식이 처음에 성립할 경우 내내 유지된다.

개별 스파이크에 대해서는 이 명제가 글자 그대로 성립하지 않는다. 여분의 입력 스파이크 하나가 뉴런을 조금 일찍 발화시키면, 1밀리초의 불응기 때문에 그다음 스파이크가 사라질 수 있다. 그러나 이 효과는 약하고 믿을 수 없어서 계산에 쓸 수 없다. 발화율에 대해서는 단조성이 엄밀하다.

결과는 인버터가 없는 여느 회로와 같다. NOT이 없다. 배타적 OR(XOR)도 없다. 두 번째 입력을 더해 출력을 끌 수 없기 때문이다. 그리고 가장 곤란한 점은 \emph{활동을 활동으로 멈출 수 없다}는 것이다. 활동을 떠받치는 것을 치울 수 있을 뿐이다.

\section{전선 두 가닥: 켜짐과 꺼짐}

막다른 길에서 나가는 방법은 몸 스스로가 알려 준다. 많은 감각 기관에서 자극은 두 벌의 뉴런으로 전달된다. 시각에서는 어떤 세포는 빛이 강해질 때, 다른 세포는 약해질 때 응답한다~\cite{Schiller1992}. 이들을 \emph{켜짐} 뉴런과 \emph{꺼짐} 뉴런이라 부르자. 네트워크는 전선 $x$ 하나 대신 쌍 $(x,\bar x)$를 받으며, 네트워크가 계산할 수 없는 `없음'은 센서가 직접 알려 준다.

전선 쌍이 있으면 NOT은 공짜이다. 전선 두 가닥을 맞바꾸기만 하면 된다. AND와 OR는 각각 뉴런 둘로 만든다.
\begin{empheq}[box=\keybox]{equation}
\begin{aligned}
x\wedge y \;&\to\; \bigl(\;x\wedge y,\;\; \bar x\vee\bar y\;\bigr),\\
x\vee y \;&\to\; \bigl(\;x\vee y,\;\; \bar x\wedge\bar y\;\bigr).
\end{aligned}
\label{eq:dualrail}
\end{empheq}
오른쪽의 연산은 모두 단조이므로 명제~\ref{prop:monotone}에 어긋나지 않는다.

비동기 회로에서 이중 레일 부호에는 더 중요한 두 번째 장점이 있다. 전선 쌍의 상태는 셋이다. `예', `아니오', 그리고 둘 다 조용할 때의 \emph{`아직 답 없음'}. 수신자는 계산이 끝났다는 것을 타이머가 아니라 신호 자체로 안다. 모든 쌍에서 정확히 한 가닥이 작동하면 끝난 것이다. 자극과 자극 사이에는 센서가 조용하고, 네트워크는 침묵으로 돌아가 다음 차례를 준비한다. 비동기 전자공학에서는 이 기법으로 소자의 지연이 어떻든 올바르게 동작하는 회로를 만든다~\cite{Sparso2001}.

\begin{keyblock}
\begin{proposition}[비동기 계산]
\label{prop:dualrail}
각 입력이 `켜짐--꺼짐' 전선 쌍으로 들어오면, \nt{GLUT} 뉴런 네트워크는 입력의 어떤 논리 함수든 계산할 수 있다. 답도 전선 쌍으로 나오며, 답이 준비되었다는 것은 모든 쌍에서 전선 한 가닥이 작동했다는 것으로 보인다. 이를 위해 클록은 필요 없다. 대가는 뉴런이 약 두 배로 드는 것이다.
\end{proposition}
\end{keyblock}

예로 배타적 OR를 보자. `두 자극 가운데 정확히 하나가 변했다'이다. 답의 정방향 전선은 $(x\wedge\bar y)\vee(\bar x\wedge y)$, 역방향 전선은 $(x\wedge y)\vee(\bar x\wedge\bar y)$이다. 뉴런 여섯 개이며, 그중 어느 것도 억제를 필요로 하지 않는다(그림~\ref{fig:dualxor}).

\begin{figure}[t]
\centering
\begin{tikzpicture}[>=Stealth,x=1cm,y=1cm,
  nrn/.style={circle,draw,very thick,fill=blue!6,minimum size=0.75cm,inner sep=0pt,font=\small},
  inp/.style={font=\small}]
\node[inp] (X)  at (0,3.3) {$x$};
\node[inp] (Xb) at (0,2.2) {$\bar x$};
\node[inp] (Y)  at (0,1.1) {$y$};
\node[inp] (Yb) at (0,0)   {$\bar y$};
\node[nrn] (A1) at (3.2,3.3) {$2$};
\node[nrn] (A2) at (3.2,2.2) {$2$};
\node[nrn] (A3) at (3.2,1.1) {$2$};
\node[nrn] (A4) at (3.2,0)   {$2$};
\node[nrn] (O)  at (7.2,2.75) {$1$};
\node[nrn] (Ob) at (7.2,0.55) {$1$};
\foreach \s/\t in {X/A1,Yb/A1,Xb/A2,Y/A2,X/A3,Y/A3,Xb/A4,Yb/A4}{\draw[->,thick,blue!60!black] (\s) -- (\t);}
\draw[->,thick,blue!60!black] (A1) -- node[pos=0.45,above,sloped,font=\scriptsize,black] {$x\wedge\bar y$} (O);
\draw[->,thick,blue!60!black] (A2) -- node[pos=0.45,below,sloped,font=\scriptsize,black] {$\bar x\wedge y$} (O);
\draw[->,thick,blue!60!black] (A3) -- node[pos=0.45,above,sloped,font=\scriptsize,black] {$x\wedge y$} (Ob);
\draw[->,thick,blue!60!black] (A4) -- node[pos=0.45,below,sloped,font=\scriptsize,black] {$\bar x\wedge\bar y$} (Ob);
\draw[->,very thick,red!70!black] (O) -- (8.6,2.75) node[right,black] {$x\oplus y$: `예'};
\draw[->,very thick,red!70!black] (Ob) -- (8.6,0.55) node[right,black] {$\overline{x\oplus y}$: `아니오'};
\node[below,font=\small,gray!40!black] at (3.2,-0.55) {AND: 문턱값 $2$};
\node[below,font=\small,gray!40!black] at (7.2,-0.55) {OR: 문턱값 $1$};
\node[below,font=\small,gray!40!black] at (0,-0.55) {전선 쌍};
\end{tikzpicture}
\caption{\nt{GLUT} 뉴런만으로 이중 레일 부호 위에 만든 배타적 OR. 원 안의 숫자는 입력 하나의 가중치를 단위로 한 문턱값이다. AND 뉴런 넷이 입력의 네 조합을 알아보고, OR 뉴런 둘이 이들을 모아 답의 전선 쌍을 만든다.}
\label{fig:dualxor}
\end{figure}

\section{시계 대신 지연}

시간을 다루는 계산에는 전선의 지연과 동시성 검출기만으로 충분하다. 가장 좋은 예는 뇌가 소리가 어디서 왔는지 알아내는 방식이다.

옆쪽에 있는 음원의 소리는 먼 귀보다 가까운 귀에 먼저 닿는다. 그 차이는 방향에 따라 달라서, 음원이 바로 앞에 있으면 0이고, 사람에서 음원이 정확히 옆에 있으면 $0.22\,\text{m}/343\,\text{m/s}\approx0.6\ms$이다. 뇌는 약 10\emph{마이크로초}의 차이를 구별한다~\cite{KlumppEady1956,Grothe2010}. 그런데 뉴런 자체는 1밀리초의 몇 분의 1보다 나은 정밀도로 발화하지 못한다. 어떻게 가능한가?

왼쪽 귀에서 오는 전선을 동시성 검출기 열을 따라 왼쪽에서 오른쪽으로, 오른쪽 귀에서 오는 전선을 오른쪽에서 왼쪽으로 깔자(그림~\ref{fig:itd}). 신호는 전선을 따라 속도 $v$로 달린다. 길이 $L$인 열의 왼쪽 끝에서 거리 $x$에 있는 검출기까지 왼쪽 신호는 시간 $x/v$, 오른쪽 신호는 $(L-x)/v$ 걸린다. 오른쪽 귀에 소리가 $\delta t$만큼 늦게 왔다면 두 신호는 다음 지점에서 동시에 만난다.
\begin{empheq}[box=\keybox]{equation}
\frac{x}{v} \;=\; \frac{L-x}{v} + \delta t
\quad\Longrightarrow\quad
x \;=\; \frac{L}{2} + \frac{v\,\delta t}{2} .
\label{eq:itd}
\end{empheq}
두 신호가 창~\eqref{eq:window} 안에 도착한 검출기만 발화한다. 발화한 검출기의 번호가 곧 음원의 방향이다. 시간 차이가 \emph{위치}로 바뀌었다.

\begin{figure}[t]
\centering
\begin{tikzpicture}[>=Stealth,
  nrn/.style={circle,draw,very thick,fill=blue!6,minimum size=0.7cm,inner sep=0pt}]
\foreach \k in {0,...,6}{\node[nrn] (D\k) at (1.4*\k,0) {\scriptsize $2$};}
\node[nrn,fill=red!15] at (D4) {\scriptsize $2$};
\draw[->,thick,blue!60!black] (-1.4,0.9) node[left,black] {\small 왼쪽 귀} -- (9.2,0.9);
\draw[->,thick,orange!85!black] (9.8,-0.9) node[right,black] {\small 오른쪽 귀} -- (-0.8,-0.9);
\foreach \k in {0,...,6}{
  \draw[->,blue!60!black] (1.4*\k,0.9) -- (D\k.north);
  \draw[->,orange!85!black] (1.4*\k,-0.9) -- (D\k.south);}
\draw[->,thick,red!70!black] (D4.north east) ++(0.1,0.05) -- ++(0.5,0.45) node[above right,font=\small,black] {발화};
\node[font=\small,gray!40!black] at (4.2,-1.5) {지연 증가 $\longrightarrow$ \hspace{2.2cm} $\longleftarrow$ 지연 증가};
\pic[scale=1.4] at (-2.3,1.75) {speaker};
\node[font=\scriptsize,align=center] at (-2.3,2.35) {왼쪽 음원};
\end{tikzpicture}
\caption{소리 방향 알아내기. 두 귀에서 온 신호가 서로를 향해 달린다. 원 하나하나는 문턱값이 $2$인 동시성 검출기이다. 오른쪽 귀에 소리가 늦게 왔으면 신호는 식~\eqref{eq:itd}에 따라 가운데보다 오른쪽에서 만난다. 네트워크의 출력은 발화한 뉴런의 번호이다. 여기서는 음원이 왼쪽에 있어 왼쪽 귀에 소리가 먼저 왔다.}
\label{fig:itd}
\end{figure}

여기에는 클록도, 억제도, 이중 레일 부호조차 없다. 네트워크는 `예'나 `아니오'로 답하지 않고 많은 뉴런 중 하나를 가리킨다. 지연선과 동시성 검출기로 된 이런 회로는 새의 청각 뇌줄기에서 실제로 작동하는 것으로 보인다~\cite{Carr1990}. 마이크로초 정밀도는 뉴런 하나하나의 정밀도가 아니라, 검출기가 많고 각자 자기 지연을 본다는 데서 나온다. 포유류에서는 이런 지연의 열이 발견되지 않았다. 포유류에서는 같은 문제를 \nt{GLYC} 뉴런이 시간을 정확히 맞춰 가하는 억제로 조율된 동시성 검출기가 푸는 것으로 보인다~\cite{Brand2002,Grothe2010}. 억제가 일치 창을 어떻게 좁히는지는 \ref{sec:gaba}장에서 \nt{GABA}를 예로 살펴본다. 글리신도 수신자의 염소 채널을 열어 같은 방식으로 억제한다.

\section{메모리}

컴퓨터에는 메모리가 필요하다. 이런 네트워크에서 메모리는 스스로를 유지하는 활동이다. 서로 강하게 연결된 \nt{GLUT} 뉴런 무리를 하나 잡아 발화율 하나 $r$로 기술하자(그림~\ref{fig:recurrentgroup}a). 평형에서 $r=f(wr+I)$이고, 여기서 $w$는 무리가 자기 자신에게 주는 피드백의 세기, $I$는 바깥에서 들어오는 유입이다. 이 방정식을 곡선 $f(wr+I)$와 직선 $r$의 교점으로 그래프를 그려 풀자(그림~\ref{fig:bistable}).

\begin{figure}[t]
\centering
% (b): tau dr/dt = -r + f(w r + I), f as in fig:bistable, w=1, I=0.6; Euler dt=0.001 tau.
\begin{tikzpicture}[>=Stealth,
  nrn/.style={circle,draw,thick,fill=blue!6,minimum size=0.5cm,inner sep=0pt}]
\begin{scope}
\node[font=\small] at (-0.5,1.55) {(a)};
\foreach \k in {1,...,6}{\node[nrn] (n\k) at ({1.4+1.0*cos(60*\k)},{1.0*sin(60*\k)}) {};}
\foreach \a/\b in {1/2,2/3,3/4,4/5,5/6,6/1,1/4,2/5,3/6}{\draw[<->,blue!60!black] (n\a) -- (n\b);}
\foreach \k in {2,3,4}{\draw[->,thick,green!45!black] ($(n\k)+(-0.95,0)$) -- (n\k);}
\node[font=\small,green!45!black] at (-0.35,-0.45) {$I$};
\node[font=\Large] at (3.1,0) {$\approx$};
\node[nrn,minimum size=0.85cm,font=\small] (R) at (4.35,-0.1) {$r$};
\draw[->,thick,blue!60!black] (R) edge[loop above,looseness=6] node[above,font=\small] {$w$} (R);
\draw[->,thick,green!45!black] (3.55,-0.95) node[left,font=\small] {$I$} -- (R);
\end{scope}
\begin{scope}[xshift=6.3cm,y=2.6cm,x=1.05cm]
\node[font=\small] at (-0.6,1.25) {(b)};
\draw[->,gray!70!black] (0,0) -- (7.3,0) node[below left,font=\small,black] {$t/\tau$};
\draw[->,gray!70!black] (0,0) -- (0,1.12) node[left,font=\small,black] {$r$};
\foreach \t in {0,2,4,6}{\draw[gray!60!black] (\t,-0.02) -- (\t,0.02); \node[below,font=\scriptsize] at (\t,0) {\t};}
\draw[densely dotted,gray!60!black] (0,0.5) -- (7,0.5) node[right,font=\scriptsize,black,align=left] {쓰기\\문턱값};
\fill[green!45!black,opacity=0.25] (1,-0.2) rectangle (2,-0.13);
\fill[gray!60!black,opacity=0.35] (1,-0.29) rectangle (1.5,-0.22);
\draw[very thick,blue!60!black] plot coordinates {(0.0,0.002) (0.1,0.002) (0.2,0.002) (0.3,0.002) (0.4,0.002) (0.5,0.002) (0.6,0.002) (0.7,0.002) (0.8,0.002) (0.9,0.002) (1.0,0.002) (1.1,0.083) (1.2,0.165) (1.3,0.242) (1.4,0.313) (1.5,0.378) (1.6,0.437) (1.7,0.491) (1.8,0.539) (1.9,0.583) (2.0,0.623) (2.1,0.643) (2.2,0.665) (2.3,0.688) (2.4,0.71) (2.5,0.732) (2.6,0.753) (2.7,0.773) (2.8,0.792) (2.9,0.81) (3.0,0.826) (3.2,0.855) (3.4,0.88) (3.6,0.9) (3.8,0.917) (4.0,0.931) (4.4,0.953) (4.8,0.967) (5.2,0.977) (5.6,0.984) (6.0,0.988) (6.5,0.992) (7.0,0.994)};
\draw[very thick,gray!60!black,densely dashed] plot coordinates {(0.0,0.002) (0.1,0.002) (0.2,0.002) (0.3,0.002) (0.4,0.002) (0.5,0.002) (0.6,0.002) (0.7,0.002) (0.8,0.002) (0.9,0.002) (1.0,0.002) (1.1,0.083) (1.2,0.165) (1.3,0.242) (1.4,0.313) (1.5,0.378) (1.6,0.358) (1.7,0.336) (1.8,0.313) (1.9,0.291) (2.0,0.269) (2.2,0.228) (2.4,0.191) (2.6,0.159) (2.8,0.133) (3.0,0.11) (3.2,0.091) (3.4,0.076) (3.6,0.063) (3.8,0.052) (4.0,0.043) (4.4,0.03) (4.8,0.021) (5.2,0.015) (5.6,0.011) (6.0,0.008) (6.5,0.006) (7.0,0.004)};
\node[font=\scriptsize,blue!60!black,anchor=west] at (3.3,0.8) {유입 $1\tau$: 기록됨};
\node[font=\scriptsize,gray!50!black,anchor=west] at (2.3,0.3) {유입 $0.5\tau$: 꺼짐};
\end{scope}
\end{tikzpicture}
\caption{자기 자신과 강하게 연결된 무리. (a)~서로를 흥분시키는 많은 \nt{GLUT} 뉴런은 세기 $w$로 자기 자신을 흥분시키는 발화율 $r$의 뉴런 하나처럼 행동한다. 바깥에서 유입 $I$가 들어온다. (b)~$w=1$이고 곡선 $f$가 그림~\ref{fig:bistable}과 같을 때 무리의 발화율을 시간에 따라 그렸다. 유입 $I=0.6$은 축 아래 막대로 표시한 시간 동안 가해진다. 유입이 $1\tau$ 동안 유지되면 무리는 불안정한 점 $r=0.5$를 넘어서 불이 붙고, 유입이 끝난 뒤에도 계속 켜져 있다. $0.5\tau$뿐이면 그 점에 이르지 못하고 다시 꺼진다. 비트를 쓰려면 유입이 약 $0.7\tau$보다 오래 지속되어야 한다.}
\label{fig:recurrentgroup}
\end{figure}

\begin{figure}[t]
\centering
\begin{tikzpicture}[>=Stealth,x=5cm,y=3.2cm]
\draw[->,gray!70!black] (0,0) -- (1.1,0) node[below left,black] {\small $r$};
\draw[->,gray!70!black] (0,0) -- (0,1.1);
\draw[thick,gray!60!black] (0,0) -- (1.05,1.05) node[right,black] {\small $r$};
\draw[very thick,blue!60!black] plot[domain=0:1.05,samples=80] (\x,{1/(1+exp(-(\x-0.5)/0.08))});
\draw[very thick,red!70!black,densely dashed] plot[domain=0:1.05,samples=80] (\x,{1/(1+exp(-(\x+0.3-0.5)/0.08))});
\fill[blue!60!black] (0.002,0.002) circle (0.018);
\draw[blue!60!black,fill=white,thick] (0.5,0.5) circle (0.018);
\fill[blue!60!black] (0.998,0.998) circle (0.018);
\node[below right,font=\small] at (0.02,0.0) {꺼짐};
\node[right,font=\small] at (0.52,0.46) {불안정};
\node[below,font=\small] at (0.99,0.96) {켜짐};
\node[anchor=east,font=\small,red!70!black] at (0.19,0.62) {$f(wr+I)$};
\node[anchor=west,font=\small,blue!60!black] at (0.45,0.1) {$f(wr)$};
\end{tikzpicture}
\caption{피드백이 강한 \nt{GLUT} 뉴런 무리로 만든 메모리. 실선 곡선은 바깥 유입이 없을 때이다. 직선 $r$과 안정한 교점 둘, 즉 `꺼짐'과 `켜짐'이 있고, 그 사이에 불안정한 교점이 있다. 짧은 유입 $I$(파선 곡선)가 가해지면 위쪽 교점만 남는다.}
\label{fig:bistable}
\end{figure}

피드백이 강하면 교점은 셋이다. 아래쪽은 침묵, 위쪽은 무리가 한껏 켜진 상태, 가운데는 불안정하다. 두 안정 상태를 가진 소자, 곧 플립플롭이 생겼다. 여기에 1을 쓰기는 쉽다. 짧은 유입이 곡선을 들어 올리면 아래쪽 교점이 사라지고, 무리에 불이 붙는다. 그리고 유입이 끝난 뒤에도 계속 켜져 있다(그림~\ref{fig:recurrentgroup}b). 너무 짧은 유입은 무리를 불안정한 점까지 끌어올리지 못하고, 무리는 다시 꺼진다. 자극이 사라진 뒤 몇 초 동안 유지되는 이런 자기 유지 활동은 실제로 뇌에서 관찰되며 단기 기억의 바탕으로 여겨진다~\cite{Wang2001}. 그러나 몇 초를 기억하는 방법이 이것만은 아니다. 기록은 시냅스 자체의 느린 변수에도 저장할 수 있다. 스파이크 버스트 뒤의 촉진은 \ref{sec:code}장의 `대시' 시냅스처럼 약 1초 동안 유지되고, 네트워크는 짧은 폭발 사이사이에 거의 침묵할 수 있다. 플립플롭이 아니라 충전된 커패시터이다~\cite{Mongillo2008}. 기억을 유지하는 동안의 이런 `조용한' 구간은 관찰되며~\cite{Stokes2015}, 끊임없는 스파이크 없이 저장하는 편이 에너지가 덜 든다. 피질이 어떤 경우에 두 방법 중 어느 것을 쓰는지는 논의 중이다.

그러면 어떻게 지우는가? 명제~\ref{prop:monotone}에 따르면 활동으로는 지울 수 없다. 무언가를 \emph{치우는} 수밖에 없다. 첫째 방법은 지지를 치우는 것이다. 무리가 다른 무리에서 오는 유입과 함께일 때만 켜져 있다면, 기억은 그 유입과 함께 꺼진다. 둘째는 피로이다. 실제 시냅스는 오래 작동하면 신경전달물질 저장량이 고갈되고~\cite{Tsodyks1997}, 뉴런에서는 흥분성을 떨어뜨리는 느린 전류가 자라난다. 피드백이 약해지고 위쪽 교점이 사라지면, 무리는 몇 초 뒤 저절로 꺼진다. 신호로 켜지고 시간으로만 꺼지는 메모리이다.

\section{순서열}

클록 대신 \emph{사건이 시동하는 타이머}를 둘 수 있다. \nt{GLUT} 뉴런 무리를 사슬로 잇자. 첫째가 둘째를, 둘째가 셋째를 흥분시키는 식이다. 바깥 신호가 첫 무리에 불을 붙이면 활동의 파동이 사슬을 따라 달린다. 각 고리는 앞 고리보다 지연 한두 개만큼 뒤에, 그리고 딱 한 번만 켜진다. 스파이크 직후 뉴런은 불응 상태이고, 파동은 이미 앞으로 나아갔기 때문이다.

이런 사슬은 사건으로부터의 시간을 잰다. $k$번째 고리가 발화했다면 시동 이후 지연 $k$개만큼의 시간이 흘렀다는 뜻이다. 그 출력을 근육에 연결하면 스스로 펼쳐지는 동작 순서열을 얻는다. 명금류가 노래할 때 뇌의 한 영역에 있는 \nt{GLUT} 뉴런들은 하나씩 엄격히 차례대로, 저마다 노래의 자기 순간에 발화한다. 이런 사슬과 매우 비슷하다~\cite{Hahnloser2002}.

클록과 달리 사슬은 시동되기 전까지 침묵하고, 한 번만 작동하며, 자신에게 연결된 것들에 대해서만 시간을 잰다. 네트워크에는 여전히 공통 박자가 없다.

\section{부족한 것}

억제 없이 \nt{GLUT} 뉴런만으로도 많은 것을 만들 수 있다. 그러나 세 가지가 부족하며, 세 가지 모두 명제~\ref{prop:monotone} 때문이다.

\emph{선택.} 네트워크는 방향 하나, 행동 하나를 골라야 한다. 두 자극이 거의 같으면 그림~\ref{fig:itd}의 네트워크에서는 두 출력이 모두 발화하고, 강한 쪽을 위해 약한 쪽을 누를 수단이 없다.

\emph{초기화.} 기억을 신호로 지울 수 없다.

\emph{안정성.} 엔지니어의 설계가 아닌 방식으로 연결이 생기는 네트워크에서는 양의 피드백 루프를 피할 수 없고, 그 안에서 활동이 사태처럼 불어날 수 있다(\ref{sec:neurontypes}장). 유일한 브레이크는 역시 피로인데, 느리고 거칠다.

세 가지 문제 모두 약은 하나이다. 스파이크를 스파이크로 끄는 뉴런이다. 그것이 \nt{GABA}이며, 이 뉴런은 계산하기 위해서가 아니라 고르고, 초기화하고, 계산을 통제하기 위해 필요하다.

\section{연습문제}

\begin{exercise}[\lvlA]
\label{ex:02:1}
\emph{소리를 위한 검출기 열.} 귀에서 온 신호가 그림~\ref{fig:itd}의 전선을 따라 속도 $5$~m/s로 달리고, 가장 큰 도착 시간 차는 $0.64\ms$이다. 검출기는 $\tau=0.5\ms$, $\theta=1.5w$인 뉴런~\eqref{eq:glutneuron}이며, 이웃한 검출기가 $10$~\textmu s를 구별하도록 배치된다. 열에는 검출기가 몇 개 있고, 소리 하나에 몇 개가 발화하는가?
\end{exercise}

\begin{exercise}[\lvlB]
\label{ex:02:2}
\emph{입력이 같지 않으면 창이 밀린다.} $\tau=0.5\ms$, $\theta=1.5$인 동시성 검출기~\eqref{eq:glutneuron}가 가중치 $1$과 $0.8$인 두 입력에서 스파이크를 하나씩 받는다. 일치 창과 그 중심을 구하라. 한쪽 귀의 입력이 다른 쪽보다 $20\,\%$ 약하면 그림~\ref{fig:itd}의 열은 몇 마이크로초 틀리는가?
\end{exercise}

\begin{exercise}[\lvlB]
\label{ex:02:3}
\emph{진동하지 않는 네트워크.} $f_i$가 감소하지 않고 $i\ne j$에서 $w_{ij}\ge0$인 네트워크 $\tau\dot r_i=-r_i+f_i\bigl(\textstyle\sum_jw_{ij}r_j+I_i\bigr)$는 단조이며 지속적으로 진동하지 않는다. 모든 순환에 억제 연결이 짝수 개 있는 네트워크는 치환 $r_i\to\pm r_i$로 이것에 귀착됨을 증명하라. 두 무리의 상호 억제는 진동할 수 있는가? \nt{GLUT}--\nt{GABA} 루프는?
\end{exercise}

\begin{exercise}[\lvlB]
\label{ex:02:4}
\emph{설계: 이중 레일 가산기.} 각 비트는 전선 쌍 $(x,\bar x)$로 전달된다. 합의 쌍 $(S,\bar S)$와 자리올림의 쌍 $(C,\bar C)$를 내는 세 비트 전가산기를 \nt{GLUT} 뉴런으로 만들고 가중치와 문턱값을 밝혀라. 뉴런 여덟 개 안에 맞추어라. 그림~\ref{fig:dualxor}처럼 AND와 OR 소자로 만들면 몇 개가 필요한가?
\end{exercise}

\begin{exercise}[\lvlB]
\label{ex:02:5}
\emph{모델링: 쓰기에 걸리는 시간.} 그림~\ref{fig:recurrentgroup}의 무리: $\tau\dot r=-r+f(r+I)$, $f(u)=1/\bigl(1+e^{-(u-0.5)/0.08}\bigr)$이고, 쓰기 전에는 아래쪽 상태에 있다. 유입 $I$를 시간 $T$ 동안 가한다. $I=0.6$에서 비트를 쓰는 가장 작은 $T$와, 그보다 작으면 어떤 $T$로도 비트가 써지지 않는 문턱 $I_c$를 수치로 구하라.
\end{exercise}

\begin{exercise}[\lvlA]
\label{ex:02:6}
\emph{타이머로서의 사슬.} 사슬의 고리는 \nt{GLUT} 뉴런 $100$개이고, 지연은 $2\ms$에 독립적인 산포 $0.2\ms$가 있다. (a)~$1$~s를 재려면 뉴런이 몇 개 필요하고 상대 오차는 얼마인가? 오차는 구간에 어떻게 의존하는가? (b)~실제 타이머의 오차는 구간과 상관없이 $5\,\%$이다. 어떤 산포 원천이 이런 결과를 주는가?
\end{exercise}

\begin{exercise}[\lvlC]
\label{ex:02:7}
\emph{탐구: 플립플롭인가 커패시터인가.} 뉴런 $100$개가 비트 하나를 $10$~s 동안 유지한다. 플립플롭: 뉴런마다 초당 스파이크 $20$개. 커패시터: 촉진 $u$가 $\tau_F=1$~s로 줄어들고, $u\ge u_{\max}/2$이면 비트가 읽히며, 갱신은 뉴런마다 스파이크 세 개짜리 버스트이다. 커패시터는 몇 배 더 경제적인가? 무리를 $200\ms$ 동안 침묵시키면 비트는 어떻게 되는가?
\end{exercise}
